\documentclass{article}
\usepackage{enumitem}
\usepackage{amsmath}
\begin{document}

\title{Chapter 9: Biomolecules}
\date{}
\maketitle

\begin{enumerate}[label=Q\arabic*.]

\item Which bond links adjacent amino acids in a polypeptide chain?
\\(A) Glycosidic bond
\\(B) Peptide bond
\\(C) Phosphodiester bond
\\(D) Hydrogen bond

\item The secondary structure of a protein refers to:
\\(A) Sequence of amino acids
\\(B) Three-dimensional folding of the polypeptide
\\(C) Regular local folding patterns like $\alpha$-helix and $\beta$-pleated sheet
\\(D) Association of multiple polypeptide subunits

\item Which of the following is NOT a reducing sugar?
\\(A) Glucose
\\(B) Fructose
\\(C) Sucrose
\\(D) Maltose

\item The enzyme that catalyzes the hydrolysis of starch into maltose is:
\\(A) Lipase
\\(B) Amylase
\\(C) Protease
\\(D) Cellulase

\item Assertion: Enzymes are biological catalysts that lower the activation energy of reactions.\\
Reason: Enzymes form enzyme-substrate complexes at the active site.
\\(A) Both true, Reason is correct explanation
\\(B) Both true, Reason is not correct explanation
\\(C) Assertion true, Reason false
\\(D) Assertion false, Reason true

\item The Km value of an enzyme represents:
\\(A) Maximum reaction velocity
\\(B) Substrate concentration at which reaction velocity is half of V$_{max}$
\\(C) Optimum pH for enzyme activity
\\(D) Number of substrate molecules converted per second

\item Competitive inhibition of an enzyme can be overcome by:
\\(A) Increasing inhibitor concentration
\\(B) Increasing substrate concentration
\\(C) Decreasing temperature
\\(D) Changing the pH

\item Which of the following correctly describes the difference between DNA and RNA?
\\(A) DNA has uracil; RNA has thymine
\\(B) DNA has deoxyribose and thymine; RNA has ribose and uracil
\\(C) DNA is single-stranded; RNA is double-stranded
\\(D) DNA has adenine; RNA does not

\item The most abundant organic molecule on Earth is:
\\(A) Protein
\\(B) DNA
\\(C) Cellulose
\\(D) Lipid

\item Phospholipids in cell membranes are amphipathic because:
\\(A) They have only hydrophilic fatty acid tails
\\(B) They have a hydrophilic phosphate head and hydrophobic fatty acid tails
\\(C) They are entirely hydrophobic
\\(D) They dissolve completely in water

\item Which amino acids are essential and must be obtained from diet in humans?
\\(A) Alanine and glycine
\\(B) Leucine, isoleucine, valine, and lysine (among others)
\\(C) Glutamic acid and aspartic acid
\\(D) Serine and cysteine

\item The Chargaff rule states that in DNA:
\\(A) A + G = T + C (purines equal pyrimidines)
\\(B) A = T and G = C
\\(C) A + T = G + C
\\(D) Both A and B above

\item Which of the following is a fibrous protein?
\\(A) Haemoglobin
\\(B) Collagen
\\(C) Insulin
\\(D) Amylase

\item Match the following biomolecules with their monomers:\\
1. Starch $\rightarrow$ (a) Glucose\\
2. Protein $\rightarrow$ (b) Amino acids\\
3. DNA $\rightarrow$ (c) Deoxyribonucleotides\\
4. Phospholipid $\rightarrow$ (d) Glycerol + fatty acids + phosphate
\\(A) 1-a, 2-b, 3-c, 4-d
\\(B) 1-b, 2-a, 3-d, 4-c
\\(C) 1-c, 2-d, 3-a, 4-b
\\(D) 1-d, 2-a, 3-b, 4-c

\item The turnover number (kcat) of an enzyme refers to:
\\(A) The number of active sites per enzyme molecule
\\(B) The number of substrate molecules converted to product per enzyme molecule per second
\\(C) The concentration of enzyme needed for half-maximal activity
\\(D) The rate of enzyme denaturation at high temperatures

\end{enumerate}
\end{document}
